\documentclass[a4paper]{article}
\usepackage{amsfonts}
\usepackage{amsmath}
\usepackage{amssymb}
\usepackage{graphicx}     

\begin{document}
  
\noindent \large Auteur: Anna Voronovska \\
\noindent \large Studierichting: Informatica\\
\noindent \large Oefeningenreeks 1D nr. 21.13\\

\medskip

\normalsize

$z^4 + 14z^2 + 44z + 25 = 4z^3$\\

$z^4 - 4z^3 + 14z^2 + 44z + 25 = 0$\\

Stap 1: Test $z = -1$:\\

$(-1)^4 - 4(-1)^3 + 14(-1)^2 + 44(-1) + 25 = 1 + 4 + 14 - 44 + 25 = 0$\\

Dus, $z = -1$ is een wortel.\\

Stap 2: Horner:\\

We delen de veelterm $z^4 - 4z^3 + 14z^2 + 44z + 25 = 0$ door $(z + 1)$ :

\begin{center}
\begin{tabular}{c|rrrr|r}
   & 1 & $-4$ & 14 & 44 & 25 \\
$-1$ &   & $-1$ & 5  & $-19$ & $-25$ \\ \hline
   & 1 & $-5$ & 19 & 25 & 0
\end{tabular}
\end{center}

De vergelijking wordt: $(z + 1)(z^3 - 5z^2 + 19z + 25) = 0$.\\

Stap 3: Herhalen voor de derdegraadsveelterm $z^3 - 5z^2 + 19z + 25 = 0$:\\

We zoeken een wortel voor $z^3 - 5z^2 + 19z + 25 = 0$:\\

Test $z = -1$ opnieuw:\\

$(-1)^3 - 5(-1)^2 + 19(-1) + 25 = -1 - 5 - 19 + 25 = 0$\\

Dus, $z = -1$ is een dubbele wortel. We passen Horner nogmaals toe:\\

\begin{center}
\begin{tabular}{c|rrr|r}
   & 1 & $-5$ & 19 & 25 \\
$-1$ &   & $-1$ & 6  & $-25$ \\ \hline
   & 1 & $-6$ & 25 & 0
\end{tabular}
\end{center}

De vergelijking is nu: $(z + 1)^2(z^2 - 6z + 25) = 0$.\\

Stap 4: Wortels zoeken:\\

Voor $z^2 - 6z + 25 = 0$ :\\

$$D = (-6)^2 - 4(1)(25) = 36 - 100 = -64$$\\

Omdat $D < 0$:\\

$\sqrt{D}=8i$\\

$z_1 = \dfrac{6 - 8i}{2} = 3 - 4i$\\

$z_2 = \dfrac{6 + 8i}{2} = 3 + 4i$\\

Voor $z+1=0$\\

$z=-1$\\

Oplossingen:\\


$ \{ -1^{(2)}, 3 + 4i, 3 - 4i \}$\\



\begin{thebibliography}{999}
%\bibitem{a} Referentie
\end{thebibliography}
\end{document} 
